\lecture{4}{10 March.}{}
\subsubsection{Running time of Kruskal's algorithm}
First, sorting all edges by weight takes \(O(m \log _2 m)\) steps. Then, iterate through all edges, if \(C_u = C_v\), do nothing, and this steps take \(O(1)\) steps. If not, then add edges and update components, which takes \(O(n)\) steps, and will run this case \(n - 1\) times. Thus, the total time:
\[
    O(m \log _2 m) + O(m) + O\left( n^2 \right) = O(m \log _2 m + n^2) = \begin{dcases}
        O(m \log m), &\text{ if } m = \Omega \left( \frac{n^2}{\log n} \right)  ;\\
        O\left( n^2 \right) , &\text{ if } m = O\left( \frac{n^2}{\log n} \right).
    \end{dcases} 
\]    
With more improvement, 
\[
    O(m \log m) \text{ for all } m. 
\]

\section{Dijkstra's Algorithm}
\subsubsection{Drawbacks of trees as networks}
\begin{itemize}
    \item [(1)] Unique path between any two vertices, which could be unnecessarily long. 
    \item [(2)] No redundancy/robustness: removing any one edge/vertex disconnects the network.
\end{itemize}

\begin{question}
    Given a network \(G = (V, E)\) and weights \(w : E \to \mathbb{R} _{\ge 0} \cup \left\{ \infty \right\} \), and given two vertices \(u, v \in V\), how do we find the cheapest path from \(u\) to \(v\)?     
\end{question}

\begin{observation}
    We may assume \(G\) is complete. If an edge \(\left\{ x, y \right\} \) is missing, add it with \(\infty \) weight.   
\end{observation}

\begin{lemma} \label{lm: cheapest path}
    If \(P\) is a cheapest path from \(u\) to \(v\), then for any \(x, y \in P\), \(P[x, y]\) is a shortest path from \(x\) to \(y\).       
\end{lemma}

\begin{algorithm}[H]
    \caption{Dijkstra's algorithm}
    \KwData{A complete graph \(G=(V, E)\).}
    \KwResult{For any vertex \(x\), the cheapest graph from \(u\) to \(x\).}
    \SetKwInput{KwIdea}{Idea}
    \KwIdea{Keep a set \(W\) of vertices which we have already found the cheapest path from \(u\). For other vertices, keep track of a tentative cheapest path, i.e. the cheapest path we have found so far.}
    \(W \gets \left\{ u \right\} \), \(t(u) \gets 0\), \(\text{prev}(u) \gets \varnothing\). 

    \For() {\(x \in V \setminus \left\{ u \right\} \) } {
        \(t(x) \gets w(u,x)\). 
        
        \(\mathrm{prev}(x) \gets u \). 
    } 

    \While(){\(W \neq V\)}{
        Let \(y = \argmin_{x \in V\setminus W} t(x) \), i.e. \(t(y)\) minimizes \(t(x)\) for all \(x \notin W\). 

        Add \(y\) to \(W\). 

        \For(){every \(x \notin W \cup \left\{ y \right\} \)} {
            \If(){\(t(x) > t(y) + w(\left\{ y,x \right\} )\)}{
                \(t(x) \gets t(y) + w(\left\{ y, x \right\} )\), \(\text{prev}(x) = y\).  
            }
        }      
    }   

    Cheapest path from \(x\) to \(u\) is \(x, \text{prev}(x), \text{prev}(\text{prev}(x)), \dots , u\) with cost \(t(x)\).       
\end{algorithm}

\begin{theorem}
    Dijkstra's algorithm terminates after \(n - 1\) rounds and returns shortest path. 
\end{theorem}

\begin{proof}
    Each round adds a new vertex to \(W\), which never gets removed, so after \(n - 1\) rounds, \(W = V\). For correctness, we can do induction on the number of rounds: 
    \begin{claim}
        At the beginning of the \(i\)-th iterartion:
        \begin{itemize}
            \item [(a)] \(\vert W \vert = i \). 
            \item [(b)] for every \(x \in V\), the vertices \(x, \text{prev}(x), \text{prev}(\text{prev}(x) ), \dots , u  \) form a path from \(x\) to \(u\) of weight \(t(x)\), where all vertices except \(x\) lie in \(W\). 
            \item [(c)] for any vertex \(z \in V \setminus W\) with \(t(z) = \min \left\{ t(x) : x \in V\setminus W \right\} \), \(t(z)\) is the minimum cost of a \(uz\)-path.
            \item [(d)] For every \(y \in V \setminus W\) and \(v \in W\), 
            \[
                t(y) \le t(v) + w(\left\{ v, y \right\}).
            \]
        \end{itemize} 
    \end{claim}   
    \begin{explanation}
        \hfill
        \begin{itemize}
            \item Base case: \(i = 1\), (a), (b), (d) are obvious. For (c), suppose not, then if \(z\) minimize \(t(\cdot)\) but there is a cheaper \(uz\)-path, say 
            \[
                u \to x \to \dots \to z \in P,
            \]
            then 
            \begin{align*}
                w(P) = w(\left\{ u, x \right\} ) + w(\left\{ x, \dots \right\} ) + \dots \ge w(\left\{ u, x \right\} ) = t(x) \ge t(z),
            \end{align*}
            which is a contradiction. 
            \item Induction step: (a) is obvious. 
            
            Now we show (b). If we didn't update \(t(x), \text{prev}(x) \) in the \((i-1)\)-th iteration, then this is true by induction. If we did update, then \(\text{prev}(x) = v_0 \) where \(v_0\) is the vertex added to \(W\) in the \((i-1)\)-th iteration, and \(t(x) = t(v_0) + w(\left\{ v_0, x \right\} ) =\) weight of path \(u \to \dots \to v_0 \to x\), and \(u\) to \(\text{prev}(v_0) \) all in \(W\) by induction, and \(v_0\) was added to \(W\) in the \((i-1)\)-th iteration. 
            
            Now we show (c). Suppose not, then \(\exists z \in V\setminus W\) minimizing \(t(z)\) but \(t(z) > d(z, u)\). Consider the cheapest path from \(u\) to \(z\), say \(P\). Let \(y\) be the first vertex in \(P\) (starting from \(u\)) that is not in \(W\). (such \(y\) exists since \(z\) is not in \(W\)) Let \(x\) be the preceding vertex on \(P\), noting that \(x \in W\). Then, 
            \begin{align*}
                t(z) &> w(P) = w(P_{[u, x]}) + w(\left\{ x,y \right\} ) + w(P_{[y, z]}) \ge w(P_{[u, x]}) + w(\left\{ x, y \right\}).
            \end{align*}                        
            By \autoref{lm: cheapest path}, \(P_{[u, x]}\) has to be a lightest path from \(u\) to \(x\). By induction (where we add \(x\) to \(W\)), \(t(x) = w(P_{[u, x]})\). Thus,
            \[
                t(z) > w(P_{[u,x]}) + w(\left\{ x,y \right\} ) = t(x) + w(\left\{ x, y \right\}) \ge t(y),
            \]   
            which is by the induction statement of (d). However, this contradicts \(z\) minimizeing \(t(v)\) for \(v \in V\setminus W\). 
            
            As for (d), after the step where we added \(v\) to \(W\), we check whether \(t(v) + w(\left\{ v, y \right\} ) < t(y)\), and update if so. Hence, after this round, \(t(v)\) will not change, and \(t(v)\) is decreasing, so it is impossible that \(t(y) > t(v) + w(\left\{ v, y \right\} )\).      
        \end{itemize}
    \end{explanation}

    This implies the correctness: for any \(x\), at the iteration when we add \(x\) to \(W\), (c) implies we have the correct path for \(x\) and we never change it subsequently.    
\end{proof}